\begin{threeparttable}
\begin{tabular}{l cc}
\toprule
\multicolumn{3}{c}{\textit{RDD2: SISFOH Welfare Index as Running Variable (Adults 65+, following Bando, Galiani and Gertler 2020)}} \\
\midrule
  & TIENE\_BILLETERA & USA\_BILLETERA \\
\midrule
\addlinespace
\multicolumn{3}{l}{\textit{RDD2\_SISFOH\_baseline}} \\
  Estimate & 0.0042 & -0.0111 \\
  SE (robust) & (0.0149) & (0.0076) \\
  95\% CI & [-0.0256, 0.0328] & [-0.0282, 0.0015] \\
  N / BW & 3{,}201 / 0.77 & 2{,}286 / 0.55 \\
  Method & rdrobust & rdrobust \\
\addlinespace
\multicolumn{3}{l}{\textit{RDD2\_SISFOH\_covariates}} \\
  Estimate & 0.0035 & -0.0118 \\
  SE (robust) & (0.0154) & (0.0075) \\
  95\% CI & [-0.0295, 0.0308] & [-0.0291, 0.0003] \\
  N / BW & 2{,}941 / 0.71 & 2{,}261 / 0.56 \\
  Method & rdrobust & rdrobust \\
\bottomrule
\end{tabular}
\begin{tablenotes}
\small
\item \textit{Notes:} Running variable is a PCA-based proxy of the SISFOH welfare index constructed from ENAHO household characteristics (walls, floors, water, sanitation, electricity, smartphone, refrigerator, TV, per-capita income, overcrowding). Cutoff defined at the estimated threshold separating extreme-poor from non-extreme-poor households among adults aged 65 and above. This design replicates the identification strategy of Bando, Galiani and Gertler (2020) applied to digital wallet adoption as outcome.
\end{tablenotes}
\end{threeparttable}